\FloatBarrier
\section{Background}
\label{sec:background}
TensorRight verifies tensor graph rewrites for tensors of arbitrary rank and
size~\cite{tensorright}. Its aggregated axes represent groups containing an
arbitrary number of named axes, allowing one rewrite specification to describe
tensors with different ranks. Grouping related axes makes it possible to reason
about a family of tensors without enumerating each dimension separately.

This representation provides the background for TEPL's symbolic axis groups.
TEPL retains the ability to describe arbitrary-length groups but represents
them as ordered sequences and tracks their positions in tensor interfaces.
The order matters for constructing concrete operation attributes: a pair of
contraction lists specifies not only which axes are reduced, but which axes
are paired and in what list order. The distinction between a logical axis,
its extent, and its position is central to the method developed in
Section~\ref{sec:method}. Axis grouping itself is prior work; the contribution
pursued here is its use with ordered, compositional correspondence to
construct reusable attribute programs.
